\clearpage\fancyhead[L]{\small POLARSTERN.}\hypertarget{native-398-part-lower}{}\label{native-398-part-lower}
\begin{center}\hrule\vspace{4pt}\fontsize{10}{12}\selectfont
V. \textsc{Lindenau} u. \textsc{Bohnenberger} Zeitschrift für Astronomie u. verwandte Wiss.\par
Band IV. Seite 119..131. Juli u. August. 1817.\par\vspace{4pt}\hrule\vspace{4pt}\end{center}
\begin{flushright}Göttingen, den 18. Juli 1817.\end{flushright}
\par\indent Wenn ich nicht irre, habe ich Ihnen, theuerster Freund, meine Beobachtungen des Polarsterns am zwölfzolligen \textsc{Reichenbach}schen Kreise noch nicht vollständig mitgetheilt: ich freue mich, dass ich \SourceLeaf{399} jetzt eine Reihe Solstitialbeobachtungen damit verbinden kann, die in diesem Jahre durch das schönste Wetter mehr als je begünstigt wurden. Diese Beobachtungen sind auf der neuen Sternwarte gemacht, deren Bau übrigens, wie Sie wissen, noch nicht ganz vollendet ist, und deren eigentliche Thätigkeit erst von der Zeit datiren wird, wo die neuen Meridian-Instrumente angelangt und aufgestellt sein werden. Die Mitte der neuen Sternwarte liegt nach meiner geodätischen Bestimmung $6''11$ südlich, und $1^{\mathrm s}90$ östlich von der Mitte der alten. Der Platz, wo ich den Polarstern beobachtete, liegt $0''12$ nördlich, der Platz der Sonnenbeobachtungen $0''16$ südlich von der Mitte.

\par\indent Sie erinnern sich, dass in einigen meiner frühern Briefe von der Muthmassung die Rede gewesen ist, der so viel besprochene Unterschied der Polhöhe aus Beobachtungen der Circumpolarsterne und aus Sonnenbeobachtungen könne bei den kleinern Vervielfältigungskreisen seinen Grund in einer Biegung der Objectivseite des Fernrohrs haben, da man am Objectiv als Gegengewicht nicht blos für das prismatische Ocular, sondern auch für die Klammer einen starken messingenen Ring vorsteckt. Falls dadurch eine nachtheilige Biegung hervorgebracht werden sollte, so müsste nach \textsc{Reichenbachs} Rathe dies vorgesteckte Gewicht so viel vermindert werden, dass dadurch blos das Fernrohr gehörig äquilibrirt würde, und die Klammer durch ein am Alhidadenkreise angeschraubtes Gewicht zu balanciren. Zu meiner Absicht war es jedoch noch zweckmässiger, jenes vorgesteckte Gewicht ganz wegzunehmen, und das ganze Gleichgewicht durch ein am Alhidadenkreise angeschraubtes Gewicht zu bewirken. Denn es ist klar, dass, wenn jene Vermuthung gegründet gewesen wäre, nunmehr durch die nicht aufgehobene Biegung der Ocularseite des Fernrohrs ein Fehler im entgegengesetzten Sinn hervorgehen müsste. Wenn vorher wegen einer Biegung alle Zenithdistanzen zu klein gemessen wären, so müssten sie nun alle zu gross werden und Polarstern-Beobachtungen eine zu kleine, die Sonne eine zu grosse Polhöhe geben. Allein der Erfolg hat dies durchaus nicht bestätigt; die auf der neuen Sternwarte aus dem Polarstern, bei angeschraubtem Gewicht, bestimmte Polhöhe ist fast haarscharf dieselbe, die aus der Uebertragung der mit angestecktem Gewicht auf der alten Sternwarte beobachteten folgt, und die Sonnenbeobachtungen weichen in demselben Sinn und um dieselbe Grösse ab wie vorher. Ich hatte anfangs, ehe das Wetter zusammenhängende astronomische Beobachtungen erlaubte, einen andern Weg gewählt, jene Muthmassung zu prüfen. Ich hatte nemlich mehrere Reihen Zenithdistanzen eines terrestrischen Gegenstandes beobachtet, einige mit vorgestecktem, die andern mit angeschraubtem Gewicht. Die ersten Reihen gaben wirklich einen Unterschied in dem Sinn und ungefähr von der Grösse wie er sein müsste, wenn jene Quelle die wahre gewesen wäre. Allein bei den nachfolgenden Reihen verschwand dieser Unterschied fast ganz, und überhaupt gaben diese Beobachtungen des irdischen Gegenstandes, auf welchem sich nicht so scharf pointiren liess wie auf Objecte am Himmel, keine so gute Harmonie unter sich, wie zu dieser Prüfung nöthig gewesen wäre.

{\rightskip=0pt plus 1em \par\indent Ist dieser äusserst merkwürdige Unterschied dem Instrument zu\-zuschreiben, so muss in seinem Bau etwas sein, was bewirkt, dass wir die Zenithdistanzen damit zu klein beobachten. Jenes Gewicht war nicht schuld; mehrere andere Umstände, die ich in Verdacht zu ziehen geneigt war, habe ich gleichfalls nach sorgfältiger Prüfung einflusslos gefunden. Ich gestehe, dass ich bis jetzt keinen gefunden habe, dem ich eine \emph{solche} Wirkung beimessen könnte. Eine Wirkung \emph{entgegengesetzter} Art würde viel leichter zu erklären sein, und kann allerdings, wenn man nicht auf seiner Hut ist, statt finden. Ich selbst habe, wie Sie wissen, im vorigen Winter eine solche Erfahrung gemacht. Der Kreis war von mir ganz zerlegt, gereinigt und wieder zusammengesetzt worden. Meine ersten Beobachtungen des Polarsterns gaben mir wiederholt eine um $20''$ zu kleine Polhöhe. Es war also etwas da, was die Zenithdistanzen vermehrte. Ich fand nach mehreren andern Prüfungen die Quelle des Fehlers bald in einer Biegung der Ocularröhre, die an ihrem äussersten Ende im Fernrohr durch die Feder fest gedrückt \SourceLeaf{400} wird. Diese mochte etwas lahm geworden sein; nachdem ihr durch einige Biegung die nöthige Spannkraft wiedergegeben war, fiel der Fehler sogleich weg. Auch wenn der Faden, an dem die Zenithdistanzen beobachtet werden, nicht gehörig gespannt ist, wird derselbe Fehler erfolgen. Der Faden hängt dann in der Mitte durch sein eignes Gewicht nieder; und man beobachtet offenbar die Zenithdistanzen zu gross. Es braucht wohl nicht erinnert zu werden, dass hier nicht von groben Fehlern die Rede ist; allein wenn der schlaffe Faden in der Mitte ein Paar Secunden niederhängt, während seine ganze Länge vielleicht über einen Grad beträgt, so erkennt das Auge noch keine Krümmung, es sei denn, dass man das Bild einer vollkommen geraden horizontalen Linie aus einer schicklichen Entfernung sich auf den Faden projiciren lässt.\par}

\par\indent Es ist aber gewiss, dass man das Instrument nicht fehlerfrei sprechen kann, ohne den entgegengesetzten Fehler andern Instrumenten aufzubürden, die einen solchen Unterschied zwischen den Stern- und Sonnenbeobachtungen gar nicht, oder von geringerer Grösse geben. Hat der 12zollige Kreis Recht, so geben die dreifüssigen Kreise mit stehender Säule und \textsc{Bessels} Kreis die Zenithdistanzen zu gross; hat der dreifüssige Kreis Recht, so gibt der zwölfzollige die Zenithdistanzen zu klein und der \textsc{Bessel}sche zu gross; hat der letzte Recht, so geben alle \textsc{Reichenbach}sche Kreise zu kleine Zenithdistanzen. Allein mit \emph{Gewissheit} kann man keinem Instrumente Recht geben, wenn man nicht entweder die mögliche Fehlerquelle an den andern Instrumenten bestimmt und unläugbar nachgewiesen, oder auch anderswoher die Unrichtigkeit der mit den andern Instrumenten gefundenen Resultate \emph{entschieden} dargelegt hat. Wie die Sachen jetzt liegen, müssen wir, däucht mir, noch gestehen, dass kein bestimmter Ausspruch zu thun ist. Die relative \emph{Wahrscheinlichkeit} wird jeder Astronom nach subjectiven Gründen vielleicht anders abwägen, denn ich finde es allerdings natürlich, dass wenn gleich gegen alle Instrumente kein bestimmter Vorwurf gerichtet werden kann, man dies am lebendigsten in Beziehung auf dasjenige fühlt, womit man selbst beobachtet und womit man also am meisten vertraut ist.

\par\indent Zwei Bemerkungen aber kann ich doch nicht umhin, hier noch beizufügen.

\par\indent \emph{Erstlich} scheint mir, dass man das Zeugniss der zwölfzolligen Kreise nicht \emph{deswegen} unbedingt verwerfen kann, \emph{weil} sie die kleinsten von diesen Instrumenten sind, wenn man nicht nachweisen kann, dass sie \emph{wegen ihrer Kleinheit} constante Fehler, die immer in demselben Sinn wirken, hervorbringen. Zufällige Fehler, die bald so bald anders wirken, könnte man eher wegen der Kleinheit erwarten, und doch zeigt die Erfahrung, dass sie eine eben so gute, oder wenigstens nicht viel schlechtere Harmonie geben, wie die grössern Instrumente. Ich weiss nicht, ob man \emph{im Allgemeinen} nicht bei grössern Instrumenten wegen ihrer Grösse wenigstens eben so sehr auf seiner Hut vor constanten Fehlern sein muss.

\par\indent \emph{Zweitens} gestehe ich, dass, ohne alle Rücksicht auf andere Umstände, ein Instrument \emph{deswegen}, \emph{weil} es gar keinen Unterschied zwischen Stern- und Sonnenbeobachtungen gibt, mir noch keine Präsumtion für sich zu haben scheint. Nicht zu gedenken, dass es doch noch nicht ganz ausgemacht ist, ob die Refraction für die Sonne und für die Fixsterne \emph{genau} dieselbe ist, so kommt hier noch ein anderer wichtiger Umstand in Betracht, worüber ich mit Ihnen, wie Sie sich erinnern, schon vor mehrern Jahren mich unterhalten habe, und den ich als eine mögliche Erklärungsart des Räthsels damals aufstellte, ohne eben sehr viel Gewicht darauf legen zu wollen. Ich wiederhole hier diesen Gegenstand, da bisher, so viel ich weiss, noch nicht öffentlich die Rede davon gewesen ist. Wir \emph{beobachten} den Mittelpunkt der Figur der Sonne; die Anziehungskraft der Sonne hingegen wird auf die Planeten proxime so wirken, als wäre sie im Schwerpunkt vereinigt, und der Schwerpunkt der Sonne, nicht der Mittelpunkt der Figur wird in der Ebene der Erdbahn liegen. Man darf nicht annehmen, dass beide in mathematischer Schärfe zusammenfallen, da ohne Zweifel der Sonnenkörper keine \emph{ganz} gleichartige Masse ist. Fiele der Schwerpunkt etwas südlich vom Mittelpunkt der Figur, so würde letzterer immer \SourceLeaf{401} eine kleine nördliche Breite haben, und in so fern man diese ignorirt, müsste aus
Sonnenbeobachtungen auf der nördlichen Erdhemisphäre immer eine zu kleine Polhöhe
folgen. So leicht und natürlich inzwischen diese Erklärungsart scheint, so ist sie
doch noch mit ihren Schwierigkeiten verbunden. Man müsste erst in eine tiefere
Untersuchung eingehen, in wie fern eine solche Zusammensetzung des Sonnenkörpers
möglich ist, wobei der Schwerpunkt um $\frac{1}{100}$ oder um $\frac{1}{1000}$ des
Sonnendurchmessers südlich vom Mittelpunkt der Figur fiele, ohne dass die Oberfläche
des Sonnenkörpers, nach den Bedingungen des Gleichgewichts, mehr von der Kugelgestalt
abwiche, als es nach den Beobachtungen der Fall ist. Auch die beobachteten
Mercurs-Durchgänge kommen hier in Betracht. Da die aus denselben abgeleiteten Breiten
des Mercurs sich auf die Voraussetzung gründen, dass der Mittelpunkt der Figur der
Sonne in der Ebene der Ecliptik selbst ist, so würden alle beobachteten Breiten zu
südlich sein. Die daraus abgeleitete Länge des aufsteigenden Knotens würde zu gross,
die des niedersteigenden zu klein ausfallen, und beide würden nicht übereinstimmen.
Die von Ihnen in der Mercurstheorie mitgetheilten Zahlen deuten in der That auf so
etwas hin, scheinen sich aber nicht mit einem so grossen Unterschiede, wie die
zwölfzolligen Kreise geben, zu vertragen. Freilich sind die beobachteten Durchgänge
durch den niedersteigenden Knoten in zu geringer Anzahl. Und da Sie über die von
Ihnen gewählten Beobachtungen und die Art, wie die Breiten daraus gefolgert sind,
kein Detail mitgetheilt haben, so lässt sich über den Grad der Genauigkeit, welchen
man jeder einzelnen Breite beilegen muss, nicht urtheilen. Es wäre aber wohl der
Mühe werth, die Mercurs-Durchgänge aus diesem Gesichtspunkte noch einmal auf das
sorgfältigste zu discutiren.

\par\indent Es würde ungemein interessant sein, wenn wir Beobachtungen der Sonne und der
südlichen Circumpolarsterne mit Repetitionskreisen aus der südlichen Erdhemisphäre
hätten. Die südliche Breite müsste aber so gross sein, dass die Refraction bei den
Circumpolarsternen keine erhebliche Ungewissheit hervorbringen könnte.

\par\indent Ich bemerke noch bei dieser Gelegenheit, dass Hr. Prof. \textsc{Schumacher} mir die
von seinem Bruder, Hrn. Capitain \textsc{Schumacher} in Friedrichsmark mit einem
\textsc{Reichenbach}schen, dem Ihrigen ganz ähnlichen, astronomischen Theodolithen
beobachtete Polhöhe mitgetheilt hat, welche aus Beobachtungen des Polarsterns
$=55^\circ\,58'\,41''3$, aus Sonnenbeobachtungen $=55^\circ\,58'\,36''6$ folgte.
Der Unterschied ist fast genau derselbe wie bei meinen Beobachtungen.

\Needspace{16\baselineskip}\par\indent\emph{Beobachtungen des Polarsterns in der untern Culmination auf der Göttinger
neuen Sternwarte.}

\begin{center}
\fontsize{9}{11}\selectfont
\setlength{\tabcolsep}{2pt}
\begin{tabular}{l|c|c|c|c|c|c}
1817 & \begin{tabular}{c}Anz. d.\\Beob.\end{tabular}
& \begin{tabular}{c}Scheinbare\\Zenithdistanz\end{tabular}
& Barometer & \begin{tabular}{c}Thermo-\\meter\end{tabular}
& \begin{tabular}{c}Re-\\fraction\end{tabular}
& \begin{tabular}{c}Refr. n.\\Bessel\end{tabular} \\
\hline
März 28 & 16 & $40^\circ\,7'\,22''40$ & $27^{\mathrm z}\,8^{\mathrm l}2$ & $-0^\circ5$ & $50''78$ & $50''25$ \\
April\ \ 1 & 22 & $24.98$ & $14.9$ & $+3.0$ & $50.88$ & \\
\ \ \ \ \ \ \ \ 2 & 24 & $23.38$ & $13.3$ & $+3.7$ & $50.49$ & \\
\ \ \ \ \ \ \ \ 8 & 28 & $26.36$ & $7\;2$ & $+4.4$ & $49.40$ & $48.87$ \\
\ \ \ \ \ \ \ 11 & 26 & $25.00$ & $9.3$ & $-1.4$ & $51.16$ & \\
\ \ \ \ \ \ \ 13 & 12 & $29.83$ & $8.5$ & $0$ & $50.68$ & \\
Mai\ \ \ 7 & 30 & $34.68$ & $10.9$ & $+7.5$ & $49.22$ & $48.67$ \\
\end{tabular}
\end{center}

\Needspace{12\baselineskip}\par\indent Hieraus folgt:

\begin{center}
\fontsize{9}{11}\selectfont
\setlength{\tabcolsep}{3pt}
\begin{tabular}{l|c|c|c}
1817 & wahre Zenithdist. & Scheinb. Polardist. & Polhöhe \\
\hline
März 28 & $40^\circ\,8'\,13''18$ & $1^\circ\,40'\,2''51$ & $51^\circ\,31'\,49''33$ \\
April\ \ 1 & $15.86$ & $3.70$ & $47.84$ \\
\ \ \ \ \ \ \ \ 2 & $13.87$ & $4.00$ & $50.13$ \\
\ \ \ \ \ \ \ \ 8 & $15.76$ & $5.79$ & $50.03$ \\
\ \ \ \ \ \ \ 11 & $16.12$ & $6.97$ & $50.51$ \\
\ \ \ \ \ \ \ 13 & $20.51$ & $7.26$ & $46.75$ \\
Mai\ \ \ 7 & $23.90$ & $13.50$ & $49.67$ \\
\end{tabular}
\end{center}

\SourceLeaf{402}
\par\indent Also im Mittel aus 158 Beobachtungen $51^\circ\,31'\,49''42$, oder auf die Mitte der Sternwarte reducirt

\begin{center}\fontsize{10}{12}\selectfont $51^\circ\,31'\,49''30$ [$-0.644\,Q$ Handschriftliche Aufzeichnung.]\end{center}

\par\indent Die Refraction ist aus \textsc{Carlinis}, die Polardistanz aus \textsc{Bessels} Tafeln entlehnt. Meine Beobachtungen des Polarsterns der untern Culmination von 1813 auf der alten Sternwarte, nach denselben Tafeln berechnet, geben mir die Polhöhe derselben $51^\circ\,31'\,55''48$, also auf die neue Sternwarte reducirt, 136 Beobachtungen,

\[ 51^\circ\,31'\,49''37 \]

\par\indent Die Veränderung des Balancirgewichts hat also gar keinen Einfluss auf die Beobachtungen gehabt.

\par\indent Die obere Culmination des Polarsterns habe ich noch nicht beobachten können. Am 19. Februar beobachtete ich 6 Zenith-Distanzen, 2 Stunden von der obern Culmination entfernt, deren Berechnung ich hier als ein Beispiel noch beifügen will, um mein Versprechen in der M.\,C. 1813 S. 484 zu erfüllen. Ich sehe mit Vergnügen aus dem letzten Stück Ihrer Zeitschrift, dass auch Hr. \textsc{Littrow} diese Beobachtungsart mit Vortheil angewandt hat.

\Needspace{11\baselineskip}\par\indent Meine Beobachtungszeiten nach der Uhr waren

\begin{center}\fontsize{10}{12}\selectfont\setlength{\tabcolsep}{3pt}\begin{tabular}{ccc}
$2^{\mathrm u}$ & $54^{\mathrm m}$ & $19^{\mathrm s}5$\\
 & $59$ & $55.0$\\
$3$ & $4$ & $57.0$\\
 & $8$ & $36.0$\\
 & $13$ & $24.0$\\
 & $15$ & $31.0$\\
\end{tabular}\end{center}

\par\indent (Die grossen Intervalle waren Folge des ungünstigen Wetters, das im ganzen Winter nur in einzelnen Augenblicken Sterne sichtbar werden liess). Der durchlaufene Bogen war $222^\circ\,23'\,53''$, Barometer $27^{\mathrm z}\,10^{\mathrm l}3$, Thermometer $+4^\circ5$, Culminationszeit des Sterns $0^{\mathrm u}\,55^{\mathrm m}\,23^{\mathrm s}59$. Mein Berechnungsverfahren will ich hier blos practisch erklären.

\par\indent Das Mittel aller Uhrzeiten ist $3^{\mathrm u}\,6^{\mathrm m}\,7^{\mathrm s}08$; die Unterschiede von demselben ohne Rücksicht auf das Zeichen

\begin{center}\fontsize{10}{12}\selectfont\setlength{\tabcolsep}{3pt}\begin{tabular}{cc}
$11^{\mathrm m}$ & $47^{\mathrm s}58$\\
$6$ & $12.08$\\
$1$ & $10.08$\\
$2$ & $48.92$\\
$7$ & $16.92$\\
$9$ & $23.92$\\
\end{tabular}\end{center}

\par\indent Mit diesen Unterschieden gehe ich in die bekannte Correctionstafel der Circummeridianhöhen ein (ich bediene mich der des Hr. v. \textsc{Zach} in der Attraction des montagnes), woraus ich erhalte

\begin{center}\fontsize{10}{12}\selectfont\setlength{\tabcolsep}{3pt}\begin{tabular}{c}
$273''01$\\
$75.50$\\
$2.68$\\
$15.56$\\
$104.11$\\
$173.42$\\
\end{tabular}\end{center}

\noindent das Mittel hieraus ist $107''38$, welchem in derselben Tafel das Argument $7^{\mathrm m}\,23^{\mathrm s}73$ entspricht. Ich schliesse hieraus (vermöge eines kleinen Kunstgriffs, dessen Grund man sich leicht selbst suppliren kann oder in den Götting. gel. Anz. 1815 S. 454 [März 23 S. bei den weiter unten folgenden Anzeigen in d.\,B.] \fancyhead[L]{\small SONNENBEOBACHTUNGEN.}\SourceLeaf{403} nachsehen mag), dass das Mittel der \emph{berechneten} Zenithdistanzen für die \emph{sechs} wirklichen Beobachtungszeiten proxime gleich sei dem Mittel der für die \emph{zwei} Uhrzeiten

\begin{center}\fontsize{9}{11}\selectfont\setlength{\tabcolsep}{3pt}\begin{tabular}{ccccccccccc}
$3^{\mathrm u}$ & $6^{\mathrm m}$ & $7^{\mathrm s}08$ & $-$ & $7^{\mathrm m}$ & $23^{\mathrm s}73$ & $=$ & $2^{\mathrm u}$ & $58^{\mathrm m}$ & $43^{\mathrm s}35$\\
$3$ & $6$ & $7{,}08$ & $+$ & $7$ & $23{,}73$ & $=$ & $3$ & $13$ & $30{,}81$\\
\end{tabular}\end{center}

\Needspace{12\baselineskip}\noindent welchen die Stundenwinkel $30^\circ\,49'\,56''9$ und $34^\circ\,31'\,48''3$ entsprechen. Die für diese Stundenwinkel mit der aus \textsc{Bessels} Tafel entlehnten Declination $88^\circ\,20'\,7''53$ und der Polhöhe $51^\circ\,31'\,47''0$ nach bekannten Formeln berechneten Zenithdistanzen sind

\begin{center}\fontsize{10}{12}\selectfont\begin{tabular}{lrrr}
& $37^\circ$ & $2'$ & $57''31$\\
& $37$ & $6$ & $32.75$\\\cline{2-4}
also das Mittel & $37$ & $4$ & $45.03$\\
\end{tabular}\end{center}
\noindent oder mit Rücksicht auf die durch $d\delta$, $d\varphi$ zu bezeichnenden Correctionen der zum Grunde gelegten Declination und Polhöhe\footnote{Die Coëfficienten sind mit dem mittlern Stundenwinkel aus einer von Hrn. \textsc{Westphal} für die hiesige Polhöhe berechneten Tafel entlehnt.}

\[ = 37^\circ\,4'\,45''03 - 1.00\,d\varphi + 0.83\,d\delta \]

\par\indent Das Mittel der beobachteten Zenithdistanz ist $37^\circ\,3'\,58''83$, oder von der Refraction $44''64$ befreiet

\[ = 37^\circ\,4'\,43''47 \]

\par\indent Hieraus folgt also $d\varphi = +1''56 + 0.83\,d\delta$ oder die Polhöhe $= 51^\circ\,31'\,48''56 + 0.83\,d\delta$. Die Beobachtungen der untern Culmination haben $51^\circ\,31'\,49''42 - d\delta$ gegeben, woraus $d\delta = +0''47$ folgen würde, und die Polhöhe $51^\circ\,31'\,48''95$. (Natürlich ist dies nur Beispiels halber hinzugefügt; um die Polhöhe von der zum Grunde gelegten Declination unabhängig zu machen, müssten die Beobachtungen in der Nähe der obern Culmination viel zahlreicher sein, und die Resultate würden dann nach der Methode der kleinsten Quadrate combinirt werden müssen). Man könnte auch eine Correction der zum Grunde gelegten geraden Aufsteigung mit einführen, allein man thut besser dies zu unterlassen, da dies Element gegenwärtig, Dank sei es Ihren und unsers \textsc{Bessels} Bemühungen, jetzt so scharf bestimmt ist.

\Needspace{22\baselineskip}\par\indent Ich lasse nunmehr meine Sonnenbeobachtungen folgen.

\begin{center}\fontsize{9}{11}\selectfont\setlength{\tabcolsep}{3pt}\begin{tabular}{c|c|c|c|c|c}
$1817$ & \shortstack{durchlaufener\\Bogen} & \shortstack{Summe der\\Reduction} & \shortstack{Anz. d.\\Beob.} & Barometer & Thermometer\\
\hline$\text{Juni}\ 13$ & $283^\circ\,11'\,9''75$ & $11'\,28''90$ & $10$ & $27^{\mathrm z}\,6^{\mathrm l}6$ & $20^\circ8$\\
$16$ & $451\;21\;38.75$ & $52\,1.70$ & $16$ & $13.1$ & $13.0$\\
$17$ & $563\;46\;30.75$ & $77\,32.11$ & $20$ & $11.0$ & $14.3$\\
$18$ & $449\;59\;10.75$ & $25\,45.01$ & $16$ & $7.7$ & $16.4$\\
$19$ & $337\;8\;25.25$ & $12\,31.98$ & $12$ & $7.7$ & $19.3$\\
$20$ & $562\;18\;1.50$ & $61\;30\;15$ & $20$ & $7.4$ & $17.0$\\
$21$ & $337\;0\;25.50$ & $19\,20.42$ & $12$ & $7.5$ & $21.0$\\
$22$ & $280\;43\;13.25$ & $7\,55.97$ & $10$ & $8.3$ & $22.0$\\
$23$ & $337\;14\;3.75$ & $26\,26.97$ & $12$ & $8.1$ & $20.7$\\
$24$ & $449\;48\;49.25$ & $31\,1.55$ & $16$ & $8.1$ & $22.5$\\
$25$ & $281\;19\;53.25$ & $18\,40.57$ & $10$ & $8.1$ & $20.4$\\
$26$ & $394\;13\;39.25$ & $23\,47.90$ & $14$ & $7.5$ & $19.0$\\
$27$ & $620\;53\;26.75$ & $73\,42.54$ & $22$ & $5.8$ & $17.5$\\
$29$ & $452\;46\;31.50$ & $40\,2.88$ & $16$ & $8.9$ & $16.0$\\
$30$ & $226\;51\;23.75$ & $21\,7.83$ & $8$ & $7.6$ & $19.0$\\
\end{tabular}\end{center}

\SourceLeaf{404}
\Needspace{23\baselineskip}\par\indent Hieraus finde ich ferner

\begin{center}\fontsize{9}{11}\selectfont\setlength{\tabcolsep}{3pt}\begin{tabular}{c|c|c|c|c|c}
$1817$ & Refraction & Parallaxe & Breite & \shortstack{Red. auf d.\\Solstitium} & \shortstack{Zenithdistanz\\im Solstitium}\\
\hline$\text{Juni}\ 13$ & $29''22$ & $4.01$ & $-0''11$ & $14'\,30''92$ & $28^\circ\,3'\,52''26$\\
$16$ & $30.73$ & $3.99$ & $+0.26$ & $5\,59.44$ & $48.63$\\
$17$ & $30.33$ & $3.99$ & $+0.34$ & $3\,58.31$ & $55.30$\\
$18$ & $29.67$ & $3.99$ & $+0.38$ & $2\,21.96$ & $53.46$\\
$19$ & $29.25$ & $3.99$ & $+0.39$ & $1\,10.43$ & $54.66$\\
$20$ & $29.17$ & $3.99$ & $+0.37$ & $0\,23.71$ & $51.41$\\
$21$ & $28.99$ & $3.99$ & $+0.30$ & $0\,1.81$ & $48.91$\\
$22$ & $28.91$ & $3.99$ & $+0.20$ & $0\,4.75$ & $52.10$\\
$23$ & $29.11$ & $3.98$ & $+0.08$ & $0\,32.50$ & $51.77$\\
$24$ & $28.89$ & $3.98$ & $-0.06$ & $1\,24.96$ & $51.62$\\
$25$ & $29.12$ & $3.99$ & $-0.22$ & $2\,42.22$ & $49.96$\\
$26$ & $29.35$ & $3.99$ & $-0.37$ & $4\,24.19$ & $51.61$\\
$27$ & $29.47$ & $3.99$ & $-0.52$ & $6\,30.75$ & $53.49$\\
$29$ & $30.05$ & $4.01$ & $-0.79$ & $11\,57.68$ & $51.86$\\
$30$ & $29.60$ & $4.01$ & $-0.88$ & $15\,17.82$ & $53.88$\\
\end{tabular}\end{center}

\par\indent Im Mittel aus 214 Beobachtungen

\[ = 28^\circ\,3'\,52''15 \]

\par\indent Meine Wahrscheinlichkeitstheorie gibt mir den wahrscheinlichen Fehler dieses Resultats $= \pm 0''32$. Nach den Mailänder Ephemeriden (woraus auch Hr. Capitain \textsc{Schumacher} bei der Reduction seiner oben erwähnten Beobachtungen die Sonnendeclinationen entlehnte) ist die scheinbare Schiefe der Ekliptik $= 23^\circ\,27'\,52''9$, also Polhöhe $51^\circ\,31'\,45''05$ oder auf den Mittelpunkt reducirt

\begin{center}\fontsize{10}{12}\selectfont $50^\circ\,31'\,45''21$ [$+0.472\,Q$ Handschriftliche Aufzeichnung.]\end{center}

\noindent um $4''09$ kleiner als aus dem Polarstern. Nach \textsc{Bessels} Schiefe der Ekliptik wäre der Unterschied noch anderthalb Secunden grösser.

\par\indent Ich muss noch bemerken, dass der Thermometerstand am Barometer nicht mit aufgezeichnet ist. Bei so mässigen Zenithdistanzen kann die Vernachlässigung der Reduction auf eine Normaltemperatur oder auf die Temperatur der Luft im Freien nur einen sehr kleinen Bruch von einer Secunde Unterschied geben. Sollte jemand die Beobachtungen nach andern Refractionstafeln und mit Rücksicht auf eine solche Refraction berechnen wollen, so wird es wenig fehlen, wenn er jedesmal bei den Sonnenbeobachtungen die Temperatur des Quecksilbers $2^\circ$ niedriger, bei den Polarsternbeobachtungen $2^\circ$ höher voraussetzt als die Temperatur im Freien. Zu dieser Bemerkung veranlasst mich die an sich sehr gegründete Erinnerung des Hrn. Prof. \textsc{Pasquich} in dem mir eben jetzt eingehändigten neuen Heft der Zeitschrift. [Handschriftlich: $51^\circ\,30'\,48''26 - 0.966\,Q$ Wintersolst. 1817.]
